% ============================================================================
% EnerTEF Journal Manuscript v2
% Merge of:
%   Paper 1 (70%) — An Autonomous Continual-Learning Framework for EV Charging
%   Paper 2 (30%) — Representation Beats Regularization in Energy Transfer
% Structure follows Paper 1. Sections VI and VII carry the new contribution.
%
% RULE OBSERVED THROUGHOUT: every number is traceable to Paper 1, Paper 2, or
% Data/results/. Nothing is invented. Missing evidence is marked
% [REQUIRED EXPERIMENT] rather than written as if it existed.
% ============================================================================
\documentclass[journal]{IEEEtran}

\usepackage{amsmath,amssymb,booktabs,graphicx,url,xcolor}
\newcommand{\params}{\boldsymbol{\theta}}
\newcommand{\ewclambda}{\lambda}
\newcommand{\nrmse}{\mathrm{nRMSE}}
\newcommand{\needexp}[1]{\textcolor{red}{[REQUIRED EXPERIMENT: #1]}}

\begin{document}

% ---------------------------------------------------------------- TITLE
% Five candidate titles (pick one, delete the rest):
% 1. When Rehearsal Beats Anchoring, and When It Does Not: Representation-Contingent
%    Continual Learning for Autonomous EV Charging Forecasts
% 2. Autonomous Continual Learning for EV Charging: Selecting the Adaptation
%    Mechanism from Evidence
% 3. Representation-Contingent Continual Learning in a Deployed EV Charging System
% 4. What Retention Measures: Calibration, Coverage and Anchoring in Autonomous
%    Energy Forecasting
% 5. Champion--Challenger Continual Learning for Deployed Energy Forecasters

\title{Autonomous Continual Learning for EV Charging:\\
Selecting the Adaptation Mechanism from Evidence}

\author{Abdelhaq Dabaghia, Pedro Rodriguez%
\thanks{Luxembourg Institute of Science and Technology (LIST), L-4362
Esch-sur-Alzette, and University of Luxembourg.}%
\thanks{A preliminary version of the deployment architecture appeared in
[conference reference]. This paper extends it with the representation analysis
of Sections~\ref{sec:repr}--\ref{sec:selection}, a retention measurement
protocol, and the champion--challenger selection mechanism.}}

\maketitle

% ---------------------------------------------------------------- ABSTRACT
\begin{abstract}
Forecasting models deployed at commercial energy sites degrade as operating
conditions drift, and continual learning is the standard remedy. Which
continual-learning mechanism to deploy, however, is normally fixed at design
time. We show that this choice cannot be made a priori, because the relative
effectiveness of parameter anchoring and experience rehearsal is not a property
of the methods but of the regime in which adaptation occurs. In a $12$-cycle
sequential deployment at a Luxembourg commercial site (five seeds, $240$
adaptation jobs), bounded replay holds backward transfer near zero
($-1.39$~kW) while an output-sensitivity regulariser forgets more than naive
fine-tuning ($+9.68$ versus $+7.04$~kW), and no regularisation strength in
$\{0,\dots,10^4\}$ restores retention. In a cross-site transfer study on the
same production models, the ordering inverts: under an instance-normalised
representation, every one of eight regularisation variants retains better than
replay (Holm-corrected $p{=}0.023$, paired $d_z$ to $-7.75$, unanimous over ten
seeds). We further show that a substantial share of reported retention
differences is not knowledge loss at all: the deployed model under-predicts its
own source holdout by $15.2$~kW at correlation $0.958$, and a pure affine
rescale recovers most of the apparent gap. We therefore propose a measurement
protocol that reports the pre-adaptation reference and separates calibration
drift from shape loss, and extend the existing tolerance-gated promotion of the
deployed pipeline into a champion--challenger mechanism that trains one
candidate per adaptation family each cycle and lets the validation gate select
among them. Finally, we bound what forecasting improvement is worth downstream:
under real ENTSO-E day-ahead prices, the gap between naive persistence and
perfect foresight is $1.74$~EUR/day, or $1.55\%$ of the bill. The results argue
that autonomous energy forecasters should select their adaptation mechanism from
online evidence rather than from assumption.
\end{abstract}

\begin{IEEEkeywords}
Continual learning, autonomous AI systems, catastrophic forgetting,
representation learning, experience replay, model predictive control,
electric vehicle charging, concept drift.
\end{IEEEkeywords}

% ---------------------------------------------------------------- I. INTRO
\section{Introduction}
\label{sec:intro}

Deployed energy forecasters are static artefacts in a non-stationary
environment. Consumption patterns drift, seasons turn, charging infrastructure
is extended, and the distribution of incoming data diverges from the one the
model was trained on. Continual learning (CL) addresses this by adapting the
model incrementally, and a substantial literature now offers mechanisms to do
so while limiting catastrophic forgetting.

Deploying such a mechanism, however, requires committing to one. A production
pipeline adapts nightly with a specific method --- a proximal regulariser, a
rehearsal buffer, or plain fine-tuning --- chosen at design time and rarely
revisited. This paper argues that the choice cannot be made correctly at design
time, because the ranking of these mechanisms is not stable across regimes.

We report two experiments on the same production models that yield opposite
orderings. In a sequential deployment over twelve monthly cycles at a commercial
site, bounded experience replay resists seasonal forgetting while an
output-sensitivity regulariser does not improve on naive fine-tuning at any
regularisation strength. In a cross-site transfer study, under an
instance-normalised output representation, the same regularisation family
retains substantially better than replay. Neither result is an artefact: both
rest on paired designs with unanimous per-seed signs.

A second observation complicates the picture further. Retention is conventionally
reported as the error of the adapted model on a past regime, relative to
competing methods. We show that a large share of the differences so obtained
reflects \emph{output calibration} rather than retained knowledge: the deployed
model under-predicts its own source holdout by $15.2$~kW while preserving the
signal shape at correlation $0.958$, and a single affine rescale recovers most
of the apparent retention gap. Any method that re-trains the output layer
captures part of this, and is credited with retention it did not perform.

These two findings motivate the contributions of this paper.
\emph{First}, a measurement protocol that makes retention claims comparable:
report the model's error on the past regime \emph{before any adaptation}, and
report a scale-invariant term alongside the error so that calibration drift is
distinguishable from loss of structure.
\emph{Second}, evidence that the anchoring-versus-rehearsal ordering is
regime-contingent, together with three candidate explanations that we test and
reject.
\emph{Third}, an extension of the deployed pipeline's existing validation gate
into a champion--challenger mechanism: rather than fixing the adaptation family,
the system trains one candidate per family each cycle and promotes on measured
evidence.
\emph{Fourth}, an end-to-end evaluation through the site's model predictive
control (MPC) scheduler, in which the marginal value of forecasting accuracy is
quantified rather than assumed.

% ---------------------------------------------------------------- II. RELATED
\section{Related Work}
\label{sec:related}

\subsection{Continual Learning for Energy Forecasting}
Regularisation-based methods constrain updates by estimated parameter
importance; rehearsal methods retain and replay past samples. Both families have
been applied to energy forecasting, with conflicting verdicts. Studies of
sequential adaptation within one domain report continual learning as essential
under abrupt shift, and replay as halving forgetting on heterogeneous streams
while tying on benign ones. Work on foundation models for time series reports
that continual fine-tuning consistently triggers forgetting. Conversely,
regularisation has been reported to improve stability for building load, and a
recent study of a $95$-entity power system finds \emph{no universal winner}:
replay wins on reconstruction stability, regularisation under concept drift.

This paper takes that last observation as its starting point and asks what
governs the contingency.

\subsection{Retention Measurement}
Most studies report retention as a comparison among adaptation methods. The
control that makes such numbers interpretable --- the model's score on the past
regime before any adaptation --- is rarely reported; a notable exception reports
a warm-up baseline and a forgetting ratio defined as the relative increase in
error on the warm-up set. We adopt and generalise that practice, and show what
it changes.

\subsection{Representation and Normalisation}
Instance-wise normalisation and domain-native reformulations are typically
treated as preprocessing choices. We find their effect on the CL conclusion to
exceed the differences between CL algorithms by an order of magnitude, which
places them inside the methodological question rather than before it.

\subsection{Research Gap}
\label{sec:gap}
Existing work --- including our own earlier deployment study --- reports one
ordering of adaptation mechanisms per experimental setting and generalises from
it. What is missing is a protocol that renders retention claims comparable
across settings, and an autonomous mechanism that selects the adaptation family
from evidence rather than from a design-time assumption.

% ---------------------------------------------------------------- III. SYSTEM
\section{System Architecture}
\label{sec:arch}

% NOTE: condense Paper 1 Sec. III by ~30%. Keep the three-loop figure, the site
% description, and the data layer. Cut the cloud-stack and framework-version
% narrative, which belongs in the engineering-challenges section.

The system operates at a commercial site in Luxembourg combining photovoltaic
generation, two EV charging installations, and an energy-management controller.
Three loops run at different time scales: a real-time control loop, a nightly
adaptation loop, and a supervisory validation loop that governs whether an
adapted model is deployed.

\subsection{Data Layer}
Fifteen-minute smart-meter data covering PV production and EV consumption at the
two installations, joined with weather and calendar features. The same feature
construction serves both the production models and the experiments reported
here.

% ---------------------------------------------------------------- IV. CL
\section{Continual Learning Framework}
\label{sec:cl}

\subsection{Problem Statement and Notation}
Let $\params_k$ denote the model parameters after adaptation cycle $k$, and
$\mathcal{D}_k$ the data window available at that cycle. Adaptation seeks
parameters that fit $\mathcal{D}_k$ without discarding behaviour learned on
earlier windows. Retention is measured on a \emph{fixed reference window}
$\mathcal{R}$ that the stream passed at cycle $0$, through backward transfer
\begin{equation}
\mathrm{BWT}(k) \;=\; \mathrm{MAE}_{\mathcal{R}}(k) - \mathrm{MAE}_{\mathcal{R}}(0),
\label{eq:bwt}
\end{equation}
positive values denoting forgetting. We stress that
$\mathrm{MAE}_{\mathcal{R}}(0)$ is a measured quantity and must be reported: it
is the reference against which any retention claim acquires meaning
(Section~\ref{sec:protocol}).

\subsection{Forecasting Models}
% Keep Paper 1's PV LSTM and EV CNN-LSTM subsections, condensed.
A PV LSTM and an EV CNN--LSTM, both taken unchanged from production.

\subsection{Output-Sensitivity Regularisation}
\label{sec:osr}
The deployed adaptation objective augments the data loss with a proximal
penalty anchored at the previous cycle's parameters,
\begin{equation}
\mathcal{L}(\params) \;=\; \mathcal{L}_{\mathrm{data}}(\params;\mathcal{D}_k)
\;+\; \frac{\ewclambda}{2}\sum_i \Omega_i\,(\theta_i - \theta^{\star}_i)^2 ,
\label{eq:osr}
\end{equation}
where $\Omega_i$ is a Memory Aware Synapses-style output-sensitivity proxy
rather than an empirical Fisher estimate. We retain the Bayesian motivation of
this form by reference only: the deployed estimator is not a Fisher information
matrix, and presenting it as one would misstate the method.

% ---------------------------------------------------------------- V. PROTOCOL
\section{What Retention Measures}
\label{sec:protocol}

Reported retention differences conflate three distinct phenomena: loss of
learned structure, drift of output calibration, and additional learning on the
past regime. We separate them.

\subsection{The Pre-Adaptation Reference}
Evaluating the transferred production model on the source holdout \emph{before
any adaptation} gives $\nrmse = 0.5479$ for the EV model and $0.6126$ for the
PV model. Zero-shot target error is $0.7964$ (EV) and $66.48$ (PV).

Two consequences follow. A full fine-tune degrades EV source accuracy by only
$1.7\%$ relative to this reference --- in one-shot transfer there is little
forgetting for any mechanism to prevent. And bounded replay scores $0.4295$,
that is $21.6\%$ \emph{better} than the model it started from: having trained on
$60$ days of source data, it has re-learned the source rather than retained it.
A comparison of replay against fine-tuning that omits the reference measures the
benefit of revisiting data, not protection against forgetting.

\subsection{Calibration versus Structure}
On the same holdout, the production model under-predicts by $15.2$~kW while
correlating with the ground truth at $0.958$: the shape of the signal is
preserved and only its amplitude is wrong. Refitting a single scale factor
($\times 1.31$) takes the error from $0.5479$ to $0.3873$. Since any adaptation
that re-trains the output layer can absorb such a rescale, a share of measured
``retention'' is calibration repair. We therefore report the correlation
$\rho(\hat{y},y)$ alongside the error, and the residual after the best affine
map $a\hat{y}+b$, as a scale-invariant retention term.

This bias exists \emph{before} adaptation, on recent source data, for a model
trained earlier --- consistent with temporal drift of the deployed artefact
rather than with anything the adaptation did.

% ---------------------------------------------------------------- VI. REPR
\section{Representation-Contingent Adaptation}
\label{sec:repr}

% This section carries Paper 2's usable 30%.

\subsection{Representation Dominates Algorithm}
Holding architecture, data, budget and protocol fixed and varying only the
output representation, warm-start transfer error on the EV task ranges from
$0.95$ (production MinMax scalers) to $0.116$ (instance normalisation) --- an
eightfold swing. Differences between EWC variants, MAS, LwF, DER++ and A-GEM
under any single representation are an order of magnitude smaller.

\subsection{Why Anchoring Fails in the Deployed Representation}
A mechanistic analysis explains the failure of parameter anchoring reported in
Section~\ref{sec:results_seq}. At the transferred parameters, the cosine between
the source importance vector and the magnitude of the target-loss gradient is
$0.82$, and approximately $62\%$ of the target-gradient energy lies in the top
decile of source-important parameters. Anchoring therefore penalises precisely
the updates the new regime requires.

\subsection{The Ordering Inverts}
Re-running the full CL taxonomy under instance normalisation, with the
adaptation implementations unchanged so that only the representation differs
($130$ runs, ten seeds), reverses the ranking. Measured against the
pre-adaptation reference of $0.2032$, all eight regularisation variants improve
retention ($-4.2\%$ to $-4.9\%$) while plain fine-tuning degrades it
($+16.5\%$); replay is intermediate ($+0.9\%$). Against plain fine-tuning, every
regularisation variant retains better with Holm-corrected $p{=}0.023$, paired
$d_z$ reaching $-7.75$, and a rank-biserial correlation of $-1.000$ --- the
advantage holds on all ten seeds without exception.

\begin{table}[t]
\centering
\caption{Retention relative to the pre-adaptation reference. The ordering of
mechanism families inverts with the representation.}
\label{tab:inversion}
\begin{tabular}{@{}lcc@{}}
\toprule
& Deployed (MinMax) & Instance-normalised \\
\midrule
Pre-adaptation reference & $0.5479$ & $0.2032$ \\
\midrule
Plain fine-tune & $+1.7\%$ & $+16.5\%$ \\
Regularisation (EWC/MAS) & $+3.6$ to $+10\%$ & $\mathbf{-4.2}$ to $\mathbf{-4.9\%}$ \\
Bounded replay & $\mathbf{-21.6\%}$ & $+0.9\%$ \\
\bottomrule
\end{tabular}
\end{table}

\subsection{Explanations Tested and Rejected}
Three candidate accounts of the inversion were tested and do not survive.
\emph{Source-model quality}: varying only the number of source-training epochs
so that pre-adaptation error spans $0.83$ to $0.21$, the retention advantage of
regularisation over fine-tuning is unchanged ($-0.040$ then $-0.039$ for the
MAS-style estimator).
\emph{Directional gradient conflict}: the cosine between source and target
gradients is positive in eleven of twelve layers (median $+0.27$), so there is
little directional conflict for a conflict-aware mechanism to exploit.
\emph{Selective freezing}: freezing the sequence encoder dominates both
baselines under the deployed representation but the effect falls from $-24.7\%$
to $-1.5\%$ relative under instance normalisation, identifying it as an artefact
of the representation rather than a property of the architecture.

A fourth account --- that the deployed model's calibration bias, rather than the
representation, drives the result --- is under test: the two arms differ in both
the representation and whether the source model is the deployed artefact or one
retrained with the same recipe. \needexp{M3: two-arm calibration control, same
representation and data, differing only in the source model.}

\subsection{On a Second Task}
The same re-run on the PV transfer, in that task's own winning representation,
does \emph{not} invert. \needexp{M2: complete the ten-seed PV arm and report it
whichever way it falls; a reversal observed on one task must be reported as
regime-contingent, not universal.}

% ---------------------------------------------------------------- VII. SELECT
\section{Autonomous Deployment with Champion--Challenger Selection}
\label{sec:selection}

\subsection{Validation-Gated Promotion}
The deployed pipeline already trains a candidate each night and promotes it only
if it does not regress beyond a tolerance $\tau$ on a recent holdout, retaining
the incumbent otherwise. Formally, with $e_{\mathrm{rec}}$ the recent-window
error and $e^{\mathrm{inc}}_{\mathrm{rec}}$ the incumbent's,
\begin{equation}
\text{promote} \iff e_{\mathrm{rec}}(\text{cand}) \le (1+\tau)\,
e^{\mathrm{inc}}_{\mathrm{rec}} , \qquad \tau = 0.05 .
\label{eq:promote}
\end{equation}
This bounds, without eliminating, the risk of deploying a worse model.

\subsection{From One Candidate to One Per Family}
Sections~\ref{sec:results_seq} and~\ref{sec:repr} establish that no adaptation
family dominates across regimes, and that the governing factor is not reliably
predictable in advance. The rational response is not a better predictor of which
family to use, but to stop predicting: the pipeline already contains a selection
mechanism, and it can select among mechanisms rather than merely accept or
reject one.

Each cycle, $K$ candidates are trained, one per family
$m \in \mathcal{M} = \{\text{naive},\ \text{OSR},\ \text{replay},\
\text{replay}{+}\text{OSR}\}$, and evaluated on both the recent window and the
retained reference window. The promoted model is
\begin{equation}
m^{\star} = \arg\min_{m \in \mathcal{M}}
\Big[\, e_{\mathrm{rec}}(m) + \alpha\, e_{\mathcal{R}}(m) \,\Big]
\quad \text{s.t. Eq.~\eqref{eq:promote}} ,
\label{eq:selection}
\end{equation}
with $\alpha$ setting the operator's stability--plasticity preference. Setting
$\alpha{=}0$ recovers the deployed behaviour; $\alpha$ large prioritises
retention.

The cost is linear in $K$ and small in absolute terms: models are compact, and
the sequential study of Section~\ref{sec:results_seq} already executed $240$
adaptation jobs on the same infrastructure.

\needexp{M1: evaluate Eq.~\eqref{eq:selection} against each fixed-family chain on
the twelve-cycle stream. If per-cycle candidate metrics were retained from the
$240$-job study, the rule can be simulated offline at no additional training
cost; otherwise one re-run is required. \textbf{Without this experiment the
mechanism is a proposal, not a result, and the paper must be framed on
Sections~\ref{sec:protocol}--\ref{sec:repr} alone.}}

% ---------------------------------------------------------------- VIII. MPC
\section{MPC Integration and the Value of Forecasting}
\label{sec:mpc}

% Condense Paper 1 Sec. V; keep the convex formulation, add the value bound.

The adapted forecasts drive a convex MPC scheduler optimising EV charging under
ENTSO-E day-ahead prices. Because the scheduler's objective is evaluated against
the \emph{forecast}, comparing objective values across forecasters is not
meaningful: an optimistic forecast lowers the planned cost without lowering the
bill. We therefore settle each plan against realised production.

Under real day-ahead prices, the realised cost of a plan built on naive
persistence exceeds that of a plan built on perfect foresight by
$1.74$~EUR/day, or $1.55\%$ of the bill; under a flat tariff the gap falls to
$0.39$~EUR/day. Price variability multiplies the value of forecasting roughly
fourfold, but the absolute ceiling remains modest: $1.74$~EUR/day is the maximum
any forecasting improvement can deliver at this site, and a real model earns a
fraction of it. We report this bound because claims that better forecasts
improve energy management are common and rarely quantified.

% ---------------------------------------------------------------- IX. SETUP
\section{Experimental Setup}
\label{sec:setup}

\subsection{Sequential Deployment Study}
Twelve cycles walking from winter 2024 to winter 2025 in $30$-day steps, each
using a $90$-day training window, chained cycle to cycle, five seeds, four
strategies: $240$ adaptation jobs. Retention measured by Eq.~\eqref{eq:bwt} on a
fixed early-winter reference window.

\subsection{Cross-Site Transfer Study}
Two real published targets. Ten seeds per condition, seed-matched Wilcoxon
tests, Holm correction across the contrast family, paired $d_z$, and equivalence
testing at margins pre-declared before the campaign.

\subsection{Baselines}
Static model; naive fine-tuning ($\ewclambda{=}0$); OSR; bounded experience
replay ($B{=}2000$, reservoir sampling, chained); and naive/seasonal persistence.

% ---------------------------------------------------------------- X. RESULTS
\section{Results}
\label{sec:results}

\subsection{Sequential Retention}
\label{sec:results_seq}

\begin{table}[t]
\centering
\caption{Multi-cycle retention and plasticity (PV LSTM, mean$\pm$s.d.\ over five
seeds). BWT in kW; positive $=$ forgetting.}
\label{tab:multicycle}
\begin{tabular}{@{}lccc@{}}
\toprule
\textbf{Strategy} & \textbf{BWT}$_{\mathrm{mid}}$ & \textbf{BWT}$_{\mathrm{peak}}$
& \textbf{Plasticity} \\
\midrule
Naive ($\ewclambda{=}0$) & $+7.04\pm7.18$ & $+13.07\pm8.52$ & $8.54$ \\
OSR ($\ewclambda{=}10^3$) & $+9.68\pm6.37$ & $+15.63\pm4.66$ & $9.37$ \\
Replay & $\mathbf{-1.39\pm7.70}$ & $\mathbf{6.83\pm6.94}$ & $13.82$ \\
Replay${+}$OSR & $\mathbf{-1.77\pm6.64}$ & $\mathbf{5.67\pm5.41}$ & $13.04$ \\
\bottomrule
\end{tabular}
\end{table}

Bounded replay holds mid-stream backward transfer near zero while naive and OSR
degrade substantially. The regulariser does not improve retention over naive
fine-tuning; if anything it is marginally worse. Replay reduces mid-stream BWT
by $11.07$~kW relative to OSR with the same sign on all five seeds (paired
$t{=}6.88$, $p{=}0.0023$).

A sweep over $\ewclambda \in \{0,10,10^2,10^3,10^4\}$ leaves mid-stream BWT
essentially flat at $+8$ to $+12$~kW: the failure is a property of parameter
anchoring in this regime, not of the operating point. A parameter-space
diagnostic supports this reading --- cumulative weight drift
$\lVert\params_k-\params_0\rVert$ grows at the same rate across all four
strategies while retention diverges --- and a DER++-style variant lands between
OSR and replay without separating from either at three seeds.

Retention is bought with plasticity: replay's final recent-window MAE is
$13.8$~kW against $8.5$~kW for naive.

\subsection{Cross-Site Transfer and the Inversion}
Reported in Section~\ref{sec:repr}, Table~\ref{tab:inversion}.

\subsection{Comparison against Persistence}
\label{sec:persistence}
On the identical holdout, metric and horizon, naive persistence achieves
$0.3245$ (EV) and $0.3500$ (PV). In the deployed representation the transferred
models do not beat it ($0.5130$ on EV); under instance normalisation the same
architecture beats it by $2.8\times$ ($0.1157$). On the PV task no
representation examined beats it, the best (clear-sky index) reaching $0.4090$.
We report this because it isolates the paper's central claim --- the accuracy
variation attributable to representation exceeds any difference between
continual-learning algorithms by an order of magnitude --- while bounding what
the models achieve in absolute terms.

\subsection{Champion--Challenger Selection}
\needexp{M1. Results section to be written once the selection rule has been
evaluated against the fixed-family chains.}

\subsection{Deployment Reliability and MPC Operation}
% Keep Paper 1's promotion statistics and MPC operating results, condensed.

% ---------------------------------------------------------------- XI. DISC
\section{Discussion}
\label{sec:discussion}

Two mechanisms compete to explain retention. \emph{Coverage} holds that what
protects a past regime is continued exposure to its input distribution;
\emph{anchoring} holds that it is proximity in parameter space. The sequential
study supports coverage: weight drift is equal across strategies while retention
diverges, and adding the proximal penalty to replay changes neither retention
nor the ranking. The transfer study under instance normalisation supports
anchoring: there, the regularisation family retains best and replay does not.

We do not think these are contradictory so much as regime-dependent, and we are
explicit that we have not identified the governing variable. Source-model
quality, directional gradient conflict, and selective freezing were each tested
and rejected as explanations (Section~\ref{sec:repr}). The calibration account
remains open pending M3.

The practical implication does not depend on resolving the question. If the
ordering of mechanisms is contingent on factors that are not reliably
predictable before adaptation, then a deployed system should not commit to one
family. It should train candidates from several and let its validation gate ---
which exists already, and which the operator already trusts to accept or reject
a model --- select among them.

% ---------------------------------------------------------------- XII. LIMITS
\section{Limitations}
\label{sec:limits}

\emph{Absolute accuracy.} In the deployed representation our models do not beat
naive persistence at the one-step horizon, and on the PV task no representation
examined does. The one-step $15$-minute horizon is where persistence is
strongest; our contribution concerns adaptation dynamics rather than absolute
accuracy, and we do not claim the latter.

\emph{Feature construction.} The production rolling means are trailing and
inclusive, so the feature vector for step $t$ contains $y_t$. Since the
predicted step's features are supplied to the model, the reported one-step
scores benefit from information unavailable at inference.
\needexp{M4: retrain with the rolling means shifted by one step and report
whether conclusions move.}

\emph{Scope.} The sequential study covers seasonal shift at a single site with
five seeds; the transfer study covers one domain pair per task. The inversion is
established on one task and does not reproduce on the second.

\emph{Prior publication.} The deployment architecture appeared in a preliminary
conference version. Its conclusion that the regulariser does not aid retention
is reproduced here and shown to be specific to the deployed representation.

% ---------------------------------------------------------------- XIII. CONCL
\section{Conclusion}
\label{sec:conclusion}

We reported two experiments on the same deployed models that order continual
learning mechanisms oppositely, and showed that a substantial part of what is
conventionally measured as retention is output calibration rather than retained
knowledge. From this we drew a measurement protocol --- report the
pre-adaptation reference, and separate scale from structure --- and a design
consequence: an autonomous forecaster should select its adaptation mechanism
from evidence gathered online, using the validation gate it already operates,
rather than from a design-time assumption that the evidence does not support.
We further bounded the downstream value of forecasting improvement at this site
to $1.55\%$ of the electricity bill, a figure we suggest should accompany claims
that better forecasts improve energy management.

\bibliographystyle{IEEEtran}
\bibliography{references}

\end{document}
